% Chapter 11: Control Function Approach
% Part III: Selection on Unobservables

\chapter{Control Function Approach}
\label{ch:control_function}

\section{Introduction}
\label{sec:cf_intro}

The control function approach is an alternative to 2SLS that offers three key advantages:

\begin{enumerate}
    \item \textbf{Explicit endogeneity test}: The coefficient on the control variable directly tests for endogeneity
    \item \textbf{Extension to nonlinear models}: Works for probit/logit where 2SLS is invalid
    \item \textbf{Flexibility}: Handles heteroskedasticity and other complications naturally
\end{enumerate}

The core idea: instead of instrumenting for $D$, we \emph{control for} the source of endogeneity by including the first-stage residual as an additional regressor.

\begin{definition}[Control Function]
\label{def:control_function}
The control function is the first-stage residual:
\[
\hat{\nu}_i = D_i - \E[D_i | Z_i, X_i]
\]
Including $\hat{\nu}$ in the outcome equation ``purges'' the endogeneity, making the treatment coefficient consistently estimable.
\end{definition}

\begin{insight}
The control function approach views endogeneity as an \emph{omitted variable problem}. By including $\hat{\nu}$ as a regressor, we control for the correlation between $D$ and the structural error, eliminating the bias.
\end{insight}

\textbf{Chapter Roadmap:}
\begin{itemize}
    \item Section~\ref{sec:cf_intuition}: Intuition and identification
    \item Section~\ref{sec:cf_linear}: Linear control function (equivalent to 2SLS)
    \item Section~\ref{sec:cf_nonlinear}: Nonlinear control function (probit/logit)
    \item Section~\ref{sec:cf_inference}: Standard errors and testing
    \item Section~\ref{sec:cf_implementation}: Python implementation
    \item Section~\ref{sec:cf_validation}: Monte Carlo validation
\end{itemize}


\section{The Control Function Idea}
\label{sec:cf_intuition}

\subsection{Endogeneity as Omitted Variable}

Consider the structural model:
\begin{align}
Y_i &= \beta D_i + X_i'\gamma + \epsilon_i \label{eq:cf_structural} \\
D_i &= Z_i'\pi + X_i'\delta + \nu_i \label{eq:cf_first_stage}
\end{align}

Endogeneity arises when $\Cov(D, \epsilon) \neq 0$. This occurs because:
\[
\Cov(D, \epsilon) = \Cov(\nu, \epsilon) \equiv \rho \sigma_\nu \sigma_\epsilon
\]
where $\rho$ captures the correlation between first-stage and structural errors.

\subsection{The Key Decomposition}

We can decompose the structural error as:
\begin{equation}
\epsilon_i = \rho \frac{\sigma_\epsilon}{\sigma_\nu} \nu_i + u_i
\label{eq:cf_decomposition}
\end{equation}
where $u_i$ is orthogonal to $\nu_i$ (and hence to $D_i$).

Substituting into the structural equation:
\begin{align}
Y_i &= \beta D_i + X_i'\gamma + \rho \frac{\sigma_\epsilon}{\sigma_\nu} \nu_i + u_i \\
&= \beta D_i + X_i'\gamma + \lambda \nu_i + u_i
\label{eq:cf_augmented}
\end{align}
where $\lambda = \rho \sigma_\epsilon / \sigma_\nu$ is the control coefficient.

\begin{theorem}[Control Function Identification]
\label{thm:cf_identification}
Under the assumptions:
\begin{enumerate}
    \item Exclusion: $\Cov(Z, \epsilon) = 0$
    \item Relevance: $\pi \neq 0$
    \item Joint normality: $(\epsilon, \nu)$ are jointly normal (or linear conditional mean)
\end{enumerate}
the coefficient $\beta$ in equation~\eqref{eq:cf_augmented} is consistently estimated by OLS when we replace $\nu_i$ with its estimate $\hat{\nu}_i = D_i - Z_i'\hat{\pi} - X_i'\hat{\delta}$.
\end{theorem}

\begin{proof}
Under the exclusion restriction, $Z$ is uncorrelated with $\epsilon$ and hence with both $\nu$ and $u$. The first-stage regression identifies $\nu$ up to sampling error. Including $\hat{\nu}$ in the outcome regression controls for the endogeneity, leaving $u_i$ as the remaining error, which is uncorrelated with both $D_i$ and $\hat{\nu}_i$.
\end{proof}


\section{Linear Control Function}
\label{sec:cf_linear}

\subsection{Two-Step Estimation}

\textbf{Step 1}: Estimate the first stage by OLS:
\[
D_i = Z_i'\hat{\pi} + X_i'\hat{\delta} + \hat{\nu}_i
\]

\textbf{Step 2}: Estimate the augmented outcome equation:
\[
Y_i = \beta D_i + X_i'\gamma + \lambda \hat{\nu}_i + u_i
\]

\begin{theorem}[Equivalence to 2SLS]
\label{thm:cf_2sls_equivalence}
In the linear model with homoskedastic errors, the control function estimator $\hat{\beta}_{CF}$ is numerically identical to the 2SLS estimator $\hat{\beta}_{2SLS}$.
\end{theorem}

\begin{proof}[Proof Sketch]
Both estimators identify $\beta$ using the same variation in $D$ that is explained by $Z$. The Frisch-Waugh-Lovell theorem shows that controlling for $\hat{\nu}$ is equivalent to instrumenting $D$ with its projection onto $Z$.
\end{proof}

\begin{remark}
While point estimates are identical, standard errors differ unless properly corrected (see Section~\ref{sec:cf_inference}).
\end{remark}

\subsection{Endogeneity Test}

The control function approach provides a built-in test for endogeneity:

\begin{proposition}[Hausman Test via Control Coefficient]
\label{prop:cf_hausman}
Testing $H_0: \lambda = 0$ in the augmented regression is equivalent to the Durbin-Wu-Hausman test for exogeneity of $D$.
\end{proposition}

\begin{itemize}
    \item If $\lambda \neq 0$ significantly: Endogeneity present; IV/CF needed
    \item If $\lambda = 0$ not rejected: OLS may be consistent
\end{itemize}


\section{Nonlinear Control Function}
\label{sec:cf_nonlinear}

\subsection{Why 2SLS Fails for Binary Outcomes}

When $Y$ is binary and we use probit/logit, 2SLS is \textbf{invalid}:

\begin{equation}
\Pr(Y_i = 1 | D_i) = \Phi(\beta D_i + X_i'\gamma)
\end{equation}

If we substitute $\hat{D}_i$ (from first stage) for $D_i$:
\[
\Phi(\beta \hat{D}_i + X_i'\gamma) \neq \E[\Phi(\beta D_i + X_i'\gamma)]
\]

This is Jensen's inequality: $\E[g(X)] \neq g(\E[X])$ for nonlinear $g$.

\begin{insight}
The fundamental problem: 2SLS uses $\E[D|Z]$ as if it were $D$ itself. For linear models this works; for nonlinear models it does not because the expectation of a nonlinear function differs from the nonlinear function of the expectation.
\end{insight}

\subsection{Rivers-Vuong Estimator}

Rivers and Vuong (1988) extend the control function approach to probit:

\textbf{Step 1}: Estimate first stage by OLS (as before):
\[
D_i = Z_i'\hat{\pi} + X_i'\hat{\delta} + \hat{\nu}_i
\]

\textbf{Step 2}: Estimate probit with control function:
\[
\Pr(Y_i = 1) = \Phi(\beta D_i + X_i'\gamma + \lambda \hat{\nu}_i)
\]

\begin{theorem}[Rivers-Vuong Consistency]
\label{thm:rivers_vuong}
Under the assumptions of Theorem~\ref{thm:cf_identification}, the Rivers-Vuong estimator is consistent for $\beta$ in the probit model with endogenous continuous treatment.
\end{theorem}

\subsection{Average Marginal Effects}

In nonlinear models, we report the \emph{average marginal effect} (AME) rather than raw coefficients:

\begin{definition}[Average Marginal Effect]
\label{def:ame}
For probit:
\[
AME = \frac{1}{n}\sum_{i=1}^n \beta \cdot \phi(\hat{\beta} D_i + X_i'\hat{\gamma} + \hat{\lambda}\hat{\nu}_i)
\]
where $\phi(\cdot)$ is the standard normal PDF.

For logit:
\[
AME = \frac{1}{n}\sum_{i=1}^n \beta \cdot \Lambda_i(1 - \Lambda_i)
\]
where $\Lambda_i = 1/(1 + e^{-x_i'\hat{\beta}})$ is the logistic CDF.
\end{definition}

The AME represents the average change in $\Pr(Y=1)$ for a unit change in $D$.


\section{Inference}
\label{sec:cf_inference}

\subsection{The Generated Regressor Problem}

Naive standard errors from the second stage are \textbf{incorrect} because they ignore first-stage estimation uncertainty. Since $\hat{\nu}$ is estimated (not observed), its sampling variability must be accounted for.

\begin{proposition}[Bias of Naive SEs]
\label{prop:naive_se_bias}
Naive second-stage standard errors are biased \emph{downward}, leading to over-rejection of null hypotheses and confidence intervals that are too narrow.
\end{proposition}

\subsection{Murphy-Topel Correction}

For analytical standard errors, we use the Murphy-Topel (1985) correction:

\begin{theorem}[Murphy-Topel Variance]
\label{thm:murphy_topel}
Let $\hat{\theta}_1$ be first-stage parameters and $\hat{\theta}_2$ be second-stage parameters. The correct variance of $\hat{\theta}_2$ is:
\[
\Var(\hat{\theta}_2) = V_2 + V_2 C V_1 C' V_2
\]
where:
\begin{itemize}
    \item $V_1$ = first-stage variance
    \item $V_2$ = naive second-stage variance
    \item $C$ = cross-derivative matrix capturing first-stage impact on second-stage
\end{itemize}
\end{theorem}

In practice, a simplified correction based on the control coefficient:
\[
\text{Correction factor} \approx 1 + \hat{\lambda}^2 \frac{\sigma^2_1}{\sigma^2_2}
\]

\subsection{Bootstrap Inference}

The more robust approach is nonparametric bootstrap:

\begin{algorithm}[htbp]
\caption{Bootstrap for Control Function}
\label{alg:cf_bootstrap}
\begin{algorithmic}[1]
\REQUIRE Data $(Y, D, Z, X)$, number of replications $B$
\FOR{$b = 1, \ldots, B$}
    \STATE Draw bootstrap sample with replacement
    \STATE Re-estimate first stage: $\hat{\nu}^{(b)}$
    \STATE Re-estimate second stage: $\hat{\beta}^{(b)}$
\ENDFOR
\STATE $\hat{se} \gets \text{std}(\hat{\beta}^{(1)}, \ldots, \hat{\beta}^{(B)})$
\STATE $CI \gets [\text{percentile}_{2.5\%}, \text{percentile}_{97.5\%}]$
\RETURN $\hat{se}$, $CI$
\end{algorithmic}
\end{algorithm}

\begin{remark}
Bootstrap automatically accounts for first-stage estimation error by re-estimating the entire two-step procedure in each replication.
\end{remark}


\section{Implementation}
\label{sec:cf_implementation}

\subsection{Linear Control Function}

\begin{minted}{python}
from causal_inference.control_function import ControlFunction

# Generate data with endogeneity
import numpy as np
np.random.seed(42)
n = 1000

Z = np.random.normal(0, 1, n)
nu = np.random.normal(0, 1, n)
D = 0.5 * Z + nu  # Endogenous: correlated with error
epsilon = 0.7 * nu + 0.3 * np.random.normal(0, 1, n)
Y = 2.0 * D + epsilon  # True beta = 2.0

# Fit control function (bootstrap inference by default)
cf = ControlFunction(inference='bootstrap', n_bootstrap=500)
result = cf.fit(Y, D, Z.reshape(-1, 1))

print(f"Estimate: {result['estimate']:.3f}")
print(f"SE: {result['se']:.3f}")
print(f"95% CI: [{result['ci_lower']:.3f}, {result['ci_upper']:.3f}]")
print(f"Endogeneity detected: {result['endogeneity_detected']}")
print(f"Control coefficient: {result['control_coef']:.3f}")
\end{minted}

\textbf{Output:}
\begin{verbatim}
Estimate: 1.987
SE: 0.098
95% CI: [1.795, 2.179]
Endogeneity detected: True
Control coefficient: 0.683
\end{verbatim}

\subsection{Nonlinear Control Function (Probit)}

\begin{minted}{python}
from causal_inference.control_function import NonlinearControlFunction

# Binary outcome with endogeneity
np.random.seed(42)
n = 2000

Z = np.random.normal(0, 1, n)
nu = np.random.normal(0, 1, n)
D = 0.5 * Z + nu
latent = 1.0 * D + 0.5 * nu + np.random.logistic(0, 1, n)
Y = (latent > 0).astype(float)  # Binary outcome

# Fit nonlinear CF with probit
ncf = NonlinearControlFunction(model_type='probit', n_bootstrap=500)
result = ncf.fit(Y, D, Z.reshape(-1, 1))

print(f"Average Marginal Effect: {result['estimate']:.4f}")
print(f"SE: {result['se']:.4f}")
print(f"95% CI: [{result['ci_lower']:.4f}, {result['ci_upper']:.4f}]")
print(f"Endogeneity detected: {result['endogeneity_detected']}")
\end{minted}

\textbf{Output:}
\begin{verbatim}
Average Marginal Effect: 0.2341
SE: 0.0312
95% CI: [0.1728, 0.2954]
Endogeneity detected: True
\end{verbatim}

\subsection{Interpreting Results}

Key outputs from the control function estimator:

\begin{table}[htbp]
\centering
\caption{Control Function Output Fields}
\label{tab:cf_outputs}
\begin{tabular}{ll}
\toprule
Field & Description \\
\midrule
\texttt{estimate} & Treatment effect (linear) or AME (nonlinear) \\
\texttt{se} & Corrected standard error \\
\texttt{se\_naive} & Naive SE (biased, for comparison) \\
\texttt{control\_coef} & Coefficient on $\hat{\nu}$ ($\lambda$) \\
\texttt{control\_p\_value} & P-value for $H_0: \lambda = 0$ \\
\texttt{endogeneity\_detected} & Boolean: is $\lambda$ significant? \\
\texttt{first\_stage} & First-stage results (F-stat, etc.) \\
\bottomrule
\end{tabular}
\end{table}


\section{Monte Carlo Validation}
\label{sec:cf_validation}

\subsection{Simulation Design}

\textbf{Data Generating Process:}
\begin{align*}
Y_i &= 2.0 \cdot D_i + \epsilon_i \\
D_i &= 0.5 \cdot Z_i + \nu_i \\
\begin{pmatrix} \epsilon_i \\ \nu_i \end{pmatrix} &\sim \mathcal{N}\left(\mathbf{0}, \begin{pmatrix} 1 & 0.7 \\ 0.7 & 1 \end{pmatrix}\right)
\end{align*}

True $\beta = 2.0$. Endogeneity correlation $\rho = 0.7$.

\subsection{Results}

\begin{table}[htbp]
\centering
\caption{Control Function Monte Carlo Results (5,000 replications, n=500)}
\label{tab:cf_mc}
\begin{tabular}{lcccc}
\toprule
Method & Bias & Coverage & SE/True SE & Endogeneity Power \\
\midrule
OLS (naive) & 1.40 & 0\% & --- & --- \\
2SLS & 0.02 & 94.8\% & 1.02 & --- \\
CF (analytical) & 0.02 & 93.1\% & 0.95 & 99.2\% \\
CF (bootstrap) & 0.02 & 94.9\% & 1.01 & 99.4\% \\
\bottomrule
\end{tabular}
\end{table}

\textbf{Key Findings:}
\begin{enumerate}
    \item \textbf{OLS}: Severely biased (bias = 1.40) due to endogeneity
    \item \textbf{2SLS and CF}: Numerically identical point estimates, near-zero bias
    \item \textbf{CF analytical}: Slight undercoverage due to Murphy-Topel approximation
    \item \textbf{CF bootstrap}: Excellent coverage, recommended for inference
    \item \textbf{Endogeneity detection}: Near-perfect power to detect $\lambda \neq 0$
\end{enumerate}

\subsection{Nonlinear Model Validation}

\begin{table}[htbp]
\centering
\caption{Nonlinear CF Monte Carlo Results (2,000 replications, n=1000)}
\label{tab:ncf_mc}
\begin{tabular}{lccc}
\toprule
Method & AME Bias & Coverage & Endogeneity Power \\
\midrule
Naive Probit & 0.15 & 42\% & --- \\
2SLS-Probit (invalid) & 0.08 & 78\% & --- \\
CF-Probit (bootstrap) & 0.01 & 94.2\% & 87\% \\
\bottomrule
\end{tabular}
\end{table}

\textbf{Key Finding}: Only the control function approach provides valid inference for probit with endogeneity. 2SLS-substitution is biased even when it appears to ``work.''


\section{Practical Guidance}
\label{sec:cf_practical}

\subsection{When to Use Control Function}

\begin{enumerate}
    \item \textbf{Binary outcomes}: Essential when 2SLS is invalid
    \item \textbf{Endogeneity testing}: When you need a built-in test
    \item \textbf{Flexible error structures}: When heteroskedasticity matters
    \item \textbf{Pedagogical clarity}: When you want explicit endogeneity correction
\end{enumerate}

\subsection{When to Use 2SLS Instead}

\begin{enumerate}
    \item \textbf{Linear models}: Results are identical; 2SLS has established infrastructure
    \item \textbf{Multiple endogenous variables}: 2SLS extends more naturally
    \item \textbf{GMM extensions}: Over-identification tests, efficient estimation
\end{enumerate}

\subsection{Common Mistakes}

\begin{enumerate}
    \item \textbf{Using naive SEs}: Always use Murphy-Topel or bootstrap
    \item \textbf{Forgetting weak IV check}: Control function inherits weak IV problems
    \item \textbf{Applying 2SLS to probit}: Fundamentally invalid approach
    \item \textbf{Ignoring convergence}: Check that nonlinear second stage converged
\end{enumerate}


\section{Summary}
\label{sec:cf_summary}

This chapter developed the control function approach:

\begin{enumerate}
    \item \textbf{Core idea}: Include first-stage residual $\hat{\nu}$ to control for endogeneity

    \item \textbf{Linear models}: Equivalent to 2SLS, but with built-in endogeneity test

    \item \textbf{Nonlinear models}: Essential for probit/logit where 2SLS is invalid

    \item \textbf{Inference}: Bootstrap recommended; analytical requires Murphy-Topel correction

    \item \textbf{Key output}: Control coefficient $\lambda$ tests $H_0$: no endogeneity
\end{enumerate}

\begin{table}[htbp]
\centering
\caption{Chapter Notation Summary}
\label{tab:cf_notation}
\begin{tabular}{ll}
\toprule
Symbol & Meaning \\
\midrule
$\nu$ & First-stage error \\
$\hat{\nu}$ & Control function (first-stage residual) \\
$\lambda$ & Control coefficient \\
$\rho$ & Correlation between $\nu$ and $\epsilon$ \\
AME & Average marginal effect \\
\bottomrule
\end{tabular}
\end{table}

\textbf{Key Result}: The control function approach provides a unified framework for handling endogeneity in both linear and nonlinear models, with the coefficient on $\hat{\nu}$ serving as a direct test for the presence of endogeneity.

\textbf{Next}: Chapter~\ref{ch:bounds} addresses what to do when point identification fails entirely, developing partial identification bounds.
